% ------------------------------------------------------------------
%  Capitolo 11 — Capire
%  Parte III
%  Ricerca di riferimento: ricerca 01 §5.4–5.5, §8; 03 §3.4
%  Voci bibliografiche principali: chi1989, bisra2018, slamecka1978, bertsch2007, pan2023, kingshepard2025, metcalfe2017, bransford1972, chase1973
%  Epigrafe: ✔ verificata sul testo (vedi ricerca 03 e registro, ricerca 00)
%  Scheda del capitolo: ricerca/06_indice-ragionato.md, capitolo 11
% ------------------------------------------------------------------
\chapter{Capire}
\label{cap:capire}

\epigrafe[Altro è ricordare, altro è sapere.]{Aliud autem est meminisse, aliud scire.}{Seneca, Epistulae ad Lucilium 33, 8}

\iniziale{Q}{uesto capitolo} \segnaposto{apertura: da scrivere — vedi la scheda nell'indice ragionato}.

\section{Ricordare e sapere}

\segnaposto{da scrivere}

\section{Chiedersi perché}

\segnaposto{da scrivere}

\section{Spiegarsi gli esempi}

\segnaposto{da scrivere}

\section{Generare prima di ricevere}

\segnaposto{da scrivere}

\section{Sbagliare per imparare}

\segnaposto{da scrivere}

\section{Imparare a memoria non è una colpa}

\segnaposto{da scrivere}

\begin{daricordare}
\segnaposto{il principio del capitolo in due righe}
\end{daricordare}

\begin{domande}
  \item \segnaposto{domanda di recupero su questo capitolo}
  \item \segnaposto{domanda di applicazione a una materia che studi}
  \item \segnaposto{domanda di ripresa su un capitolo precedente}
\end{domande}

\begin{sintesi}
  \item \segnaposto{punto di sintesi}
\end{sintesi}

\finecapitolo
